\documentclass[variants={default,lyric},paper={a5paper},maintitle={Source Effects}]{ulsbs-songbook}
\providecommand{\SourceEffectsLyrics}{false}
\newcount\SourceEffectCount
\newcommand{\SourceEffect}{\global\advance\SourceEffectCount by 1\relax}
\newcommand{\SourceEffectsSong}{%
  \SourceEffect
  \typeout{SOURCE-SONG:\theulSongNumber}%
}
% This fixture mode suppresses all inline notation, not merely chord symbols.
\ExplSyntaxOn
\str_if_eq:eeT {\SourceEffectsLyrics}{true}
  {\ulShowChords{false}\ulShowNotes{false}\ulShowBeats{false}
   \ulShowKey{false}\ulShowGoodKeys{false}}
\ExplSyntaxOff
\begin{document}
\ulSongColumns{1}
\begin{ulSongs}{}
\ulSongChapter{First Chapter}
\beginsong{}[id={source-effects-empty},key={C}]
  \SourceEffectsSong
  \beginverse
    \[C]EmptyTitleAnchor \[G]with a chordless tail.
  \endverse
\endsong
\beginsong{A}[id={source-effects-short},key={F#m},by={Fixture Author},
  cr={Fixture Copyright},ph={II},gk={Am, B&}]
  \SourceEffectsSong
  \beginverse
    \[C]\[G]AdjacentAnchor \[F#m/C#]slash \[B&]flat \[C#]sharp.
    \[C]sylla\[G]ble hy\[Am]phen-\[E7]ation \[F]{grouped lyric words}.
    \[C\kern.5em]SpaceAnchor \[C\textsuperscript7]suffix
    \[F\shrp{}m/C\shrp{}]grouped \[F#\symdim7]diminished.
  \endverse
  \beginverse*[2]
    IndentedStarAnchor \[Am]explicit\ \[G]space and \[C]a chordless segment.
    A source newline starts this line; a percent keeps these words joined:%
    JoinedAnchor and escaped \& \% \# remain literal.
  \endverse
\endsong
\beginsong{A deliberately long title with many words and a final phrase}[id={source-effects-long},key={Am}]
  \SourceEffectsSong
  \newchords{sourceeffectsmemory}
  \beginverse\memorize[sourceeffectsmemory]
    |\[Am]StoredAnchor |\[E7]second |\[Am]third.
  \endverse
  \transpose{2}
  \beginverse\replay[sourceeffectsmemory]
    |^ReplayAnchor |^second |^third.
  \endverse
  \beginverse[1]
    \[C]WrappingAnchor this deliberately extended lyric line contains enough words to wrap at the song width while preserving chord alignment over the middle \[G]syllables and the final words.
  \endverse
  \musicnote{MusicNoteAnchor with inline \ac{1}{C} and a musical instruction.}
  \beginverse
    \altchords{|C |\ac{2}{G} |\notrans{Am}} AlternateAnchor
    \beginrep
      \[C]OuterRepeatAnchor
      \beginrep \[G]InnerRepeatAnchor \endrep
      \[Am]RepeatTailAnchor
    \endrep
  \endverse
  \brk
  \beginverse*
    AfterBreakAnchor \[D]on the next song segment.
  \endverse
\endsong
\end{ulSongs}
\begin{ulSongs}{}
\ulSongChapter{Second Chapter}
\beginsong{Cora\c{c}\~ao, luz e ora\c{c}\~ao}[id={source-effects-accented},key={D&},by={Jos\'e da Silva}]
  \SourceEffectsSong
  \beginverse+
    \[<* C D;E F @ .3em>Am]MelodyAnchor \[<A _ C;E _ G>]rest slots.
    \[<!C>]ZeroAdvanceAnchor \[</>Am]unpitched \[<* @ 1em>]beat only.
    \[<C>]\[<D>]AdjacentMelodyAnchor
  \endverse
  {\ulMelodyStyleMapping{swapped}\ulMelodySecondaryPosition{below}
    \beginverse*+[2]
      \[<C# D&;A C>Am7]SwappedAnchor \[<C @ -.4em>]negative offset.
    \endverse}
  \notesoff
  \beginverse+[1]
    \[<* C>G]HiddenNotesAnchor with a visible beat.
  \endverse
  \noteson\beatsoff
  \beginverse+
    \[<* C>G]HiddenBeatsAnchor with a visible note.
  \endverse
  \notesoff
  \beginverse*+
    \[<* C>G]HiddenBothAnchor keeps only its chord.
  \endverse
  \noteson\beatson
\endsong
\beginsong{Suomi, Portugu\^es, Espa\~nol, Fran\c{c}ais}[id={source-effects-languages},key={G}]
  \SourceEffectsSong
  \beginverse
    \[G]LanguageAnchor hyv\"a\"a y\"ot\"a, cora\c{c}\~ao, ma\~nana, lumi\`ere.
    {\SourceEffect GroupedEffectAnchor} preserves source evaluation.
  \endverse
\endsong
\beginsong{Standard Title}[id={source-effects-standard},key={C}]
  \SourceEffectsSong
  \beginverse
    \[C]FinalAnchor \[G]ordinary numbered verse.
  \endverse
  \beginverse*
    UnnumberedAnchor without a chord.
  \endverse
\endsong
\end{ulSongs}
\typeout{SOURCE-EFFECTS:\the\SourceEffectCount}
\end{document}
